% Exercises of Week 9, shared by beamer.tex, the study guide, the script and
% exercises.tex. Every exercise lives in a file of its own,
% exercises/NN-slug.tex, numbered 9.NN (see exercises.lua for the docstring).
% The body of an exercise file:
%
%   \begin{problem} ... \end{problem}      the statement, used everywhere;
%                                          inside it, \begin{context} ...
%                                          \end{context} is background that
%                                          the slides leave out
%   \begin{slideproblem} ... \end{slideproblem}
%                                          (optional) a shorter statement
%                                          for the slides
%   \begin{slides} ... \end{slides}        one frame of the deck: \exstatement
%                                          places the statement; repeat the
%                                          environment for more frames
%   \begin{solution} ... \end{solution}    the detailed solution (study
%                                          guide, exercise PDFs, website)
%
% Commands for the documents (the argument is the slug: 01-license-plates.tex
% has the slug license-plates):
%
%   \exref{slug}                 the number, e.g. 9.1, for "Exercise~\exref{...}"
%   \extitle{slug}               the title from the docstring
%   \exerciseframe[opts]{slug}   beamer: the frames of the exercise
%   \studyexercise{slug}         study guide: heading, statement, solution
%   \scriptexercise{slug}        script: the statement box, as on the slide
%   \allexercises                exercises.tex: every exercise, in order; on
%                                the website statement and solution with
%                                markers where the pipeline splits the HTML
%                                (\collectionexercise), in the PDF plain:
%                                each statement, followed by its solution if
%                                \ifsolutions (exercises-solutions.tex)

\def\exweek{9}

\directlua{exercises = dofile("exercises/exercises.lua")
  exercises.scan("exercises", "\exweek")}

\ExplSyntaxOn

\msg_new:nnn { exercises } { unknown } { Unknown~exercise~`#1'. }

\cs_new_protected:Npn \__ex_missing:n #1
  { \msg_error:nnn { exercises } { unknown } {#1} ?? }
\cs_new:Npn \exref #1
  {
    \cs_if_exist:cTF { exnum@#1 }
      { \use:c { exnum@#1 } } { \__ex_missing:n {#1} }
  }
\cs_new:Npn \extitle #1
  {
    \cs_if_exist:cTF { extitle@#1 }
      { \use:c { extitle@#1 } } { \__ex_missing:n {#1} }
  }

% ---------------------------------------------------------------------------
% Loading an exercise: its parts are stored, nothing is typeset. The slides
% typeset the stored parts; the other documents read the part they need from
% the file again (\__ex_input_part:nn), as stored tables do not survive lwarp.

\tl_new:N \g__ex_current_tl
\tl_new:N \g__ex_problem_tl
\tl_new:N \g__ex_slideproblem_tl
\tl_new:N \g__ex_solution_tl
\seq_new:N \g__ex_slides_seq
\box_new:N \l__ex_scratch_box

\cs_new_protected:Npn \__ex_load:n #1
  {
    \cs_if_exist:cF { exfile@#1 } { \__ex_missing:n {#1} }
    \tl_gset:Nn \g__ex_current_tl {#1}
    \tl_gclear:N \g__ex_problem_tl
    \tl_gclear:N \g__ex_slideproblem_tl
    \tl_gclear:N \g__ex_solution_tl
    \seq_gclear:N \g__ex_slides_seq
    \vbox_set:Nn \l__ex_scratch_box
      {
        \DeclareDocumentEnvironment { problem } { +b }
          { \tl_gset:Nn \g__ex_problem_tl {##1} } { }
        \DeclareDocumentEnvironment { slideproblem } { +b }
          { \tl_gset:Nn \g__ex_slideproblem_tl {##1} } { }
        \DeclareDocumentEnvironment { slides } { +b }
          { \seq_gput_right:Nn \g__ex_slides_seq {##1} } { }
        \DeclareDocumentEnvironment { solution } { +b }
          { \tl_gset:Nn \g__ex_solution_tl {##1} } { }
        \exp_args:Nv \input { exfile@#1 }
      }
  }

% Typesets one part (problem, slideproblem, solution) of an exercise straight
% from its file; the other parts are skipped.
\cs_new_protected:Npn \__ex_input_part:nn #1#2
  {
    \group_begin:
      \clist_map_inline:nn { problem, slideproblem, slides, solution }
        {
          \str_if_eq:nnTF {##1} {#2}
            { \DeclareDocumentEnvironment {##1} { } { } { } }
            { \DeclareDocumentEnvironment {##1} { +b } { } { } }
        }
      \exp_args:Nv \input { exfile@#1 }
    \group_end:
  }

% Exercise 9.1, or Exercise 9.2 · Finance
\cs_new:Npn \__ex_heading:
  {
    Exercise~\exref { \g__ex_current_tl }
    \cs_if_exist:cT { exfinance@ \g__ex_current_tl } { ~·~Finance }
  }
\prg_new_conditional:Npnn \__ex_if_finance: { T, F, TF }
  {
    \cs_if_exist:cTF { exfinance@ \g__ex_current_tl }
      { \prg_return_true: } { \prg_return_false: }
  }

% ---------------------------------------------------------------------------
% Beamer

% The finance context is printed in the study guide, the exercise PDFs and on
% the website; the
% slides and the script (which shows what the slide shows) leave it out.
\bool_new:N \l__ex_context_bool
\bool_set_true:N \l__ex_context_bool
\DeclareDocumentEnvironment { context } { }
  {
    \bool_if:NTF \l__ex_context_bool
      { \par \textit { Finance~context. } ~ }
      { \__ex_gobble_context:w }
  }
  { \bool_if:NT \l__ex_context_bool { \par } }
\cs_new_protected:Npn \__ex_gobble_context:w #1 \end #2
  { \str_if_eq:nnTF {#2} { context } { \end { context } } { \__ex_gobble_context:w } }

\IfClassLoadedTF { beamer }
  {
    \bool_set_false:N \l__ex_context_bool

    \cs_new:Npn \__ex_block: { \__ex_if_finance:TF { exampleblock } { block } }
    \NewDocumentCommand \exstatement { }
      {
        \begin { \__ex_block: } { \__ex_heading: }
          \tl_if_empty:NTF \g__ex_slideproblem_tl
            { \g__ex_problem_tl } { \g__ex_slideproblem_tl }
        \end { \__ex_block: }
      }

    \tl_new:N \l__ex_frames_tl
    \NewDocumentCommand \exerciseframe { O{t} m }
      {
        \__ex_load:n {#2}
        \tl_clear:N \l__ex_frames_tl
        \seq_map_indexed_inline:Nn \g__ex_slides_seq
          {
            \tl_put_right:Nn \l__ex_frames_tl { \begin{frame}[#1] }
            \tl_put_right:Ne \l__ex_frames_tl
              {
                {
                  \exp_not:N \extitle {#2}
                  \int_compare:nNnT {##1} > { 1 } { \c_space_tl (cont.) }
                }
              }
            \tl_put_right:Nn \l__ex_frames_tl { ##2 \end{frame} }
          }
        \tl_use:N \l__ex_frames_tl
      }
  }
  {
    \cs_new:Npn \__ex_box: { \__ex_if_finance:TF { fin } { task } }
    \cs_new_protected:Npn \__ex_statement_box:n #1
      {
        \begin { \__ex_box: } { \__ex_heading: }
          \__ex_input_part:nn { \g__ex_current_tl } {#1}
        \end { \__ex_box: }
      }

    \NewDocumentCommand \studyexercise { m }
      {
        \__ex_load:n {#1}
        \subsection* { Exercise~\exref{#1}~---~\extitle{#1} }
        \addcontentsline { toc } { subsection }
          { Exercise~\exref{#1}~---~\extitle{#1} }
        \__ex_statement_box:n { problem }
        \soltitle \__ex_input_part:nn {#1} { solution }
      }

    \NewDocumentCommand \scriptexercise { m }
      {
        \__ex_load:n {#1}
        \group_begin:
          \bool_set_false:N \l__ex_context_bool
          \tl_if_empty:NTF \g__ex_slideproblem_tl
            { \__ex_statement_box:n { problem } }
            { \__ex_statement_box:n { slideproblem } }
        \group_end:
      }

    % The website (lwarp): statement and solution of each exercise, between
    % markers where the pipeline splits the HTML.
    \NewDocumentCommand \collectionexercise { m }
      {
        \__ex_load:n {#1}
        \par [[exercise:#1:problem]] \par
        \__ex_input_part:nn {#1} { problem }
        \par [[exercise:#1:solution]] \par
        \__ex_input_part:nn {#1} { solution }
        \par [[exercise:#1:end]] \par
      }
    % The PDF: plain, each statement, followed by its solution in the
    % version with solutions (\ifsolutions). The statement is set off like a
    % quote in Markdown, by a bar on the left (quote=<color>, also used by
    % exercises.tex) in the color of its box in the study guide: navy, or teal
    % for a finance exercise.
    \tcbset
      {
        quote/.style =
          {
            enhanced, breakable, frame~hidden, interior~hidden,
            sharp~corners, boxrule = 0pt, borderline~west = {2pt}{0pt}{#1},
            left = 8pt, right = 0pt, top = 1pt, bottom = 1pt,
            before~skip = 6pt, after~skip = 6pt,
            % the document's paragraph spacing (parskip=half)
            before~upper = { \setlength { \parskip } { .5\baselineskip } }
          }
      }
    \NewDocumentCommand \plainexercise { m }
      {
        \__ex_load:n {#1}
        \subsection* { Exercise~\exref{#1}~---~\extitle{#1} }
        \addcontentsline { toc } { subsection }
          { Exercise~\exref{#1}~---~\extitle{#1} }
        \__ex_if_finance:TF
          { \begin{tcolorbox}[quote=teal] }
          { \begin{tcolorbox}[quote=navy] }
          \__ex_input_part:nn {#1} { problem }
        \end{tcolorbox}
        \legacy_if:nT { solutions }
          {
            \soltitle
            \__ex_input_part:nn {#1} { solution }
          }
      }
    \NewDocumentCommand \allexercises { }
      {
        \IfPackageLoadedTF { lwarp }
          { \directlua { exercises.each("collectionexercise") } }
          { \directlua { exercises.each("plainexercise") } }
      }
  }

\ExplSyntaxOff
